\documentclass[11pt]{article}
\usepackage[margin=1in]{geometry}
\usepackage{booktabs}
\usepackage{amsmath}
\usepackage{hyperref}
\title{Five Languages, Five Bottlenecks:\\ Diacritization and G2P for Transliteration at Interscript}
\author{Interscript ML Team}
\date{August 2026}

\begin{document}
\maketitle

\begin{abstract}
We modernized the phonological layer of the Interscript transliteration
platform across five languages: Arabic, Hebrew, Thai, Persian, and Urdu,
covering both diacritization (vowel restoration in the original script)
and G2P (romanization/phonemization). Our central finding is a
\emph{bottleneck taxonomy}: each language's performance ceiling was set
by exactly one resource---data \emph{quality} (Arabic), data \emph{domain}
(Hebrew), labeled-data \emph{coverage} (Thai), output
\emph{representation} (Persian, Urdu)---and a fixed ByT5-based framework
with one surgical fix per language reached or beat the published
state of the art in each. Results: Arabic \textbf{2.81\% DER} (CE) /
1.69\% (w/o CE) on the SadeedDiac-25 benchmark with Misraj's own
evaluator under a zero-skip windowed protocol---2.6$\times$/3.1$\times$
better than the 1.5B-parameter Sadeed model, ahead of GPT-4 and
Gemini-Flash-2.0 (only Claude-3.7 ahead), plus a contamination audit
showing the benchmark's public source corpus leaks 122 of its own
paragraphs; Hebrew 17.5\% DER
on Biblical text where the SOTA model scores 35.6\%; Thai 2.32\% PER
via deterministic (non-LLM) augmentation; Persian homograph accuracy
77.3--89.5\% and diacritization 0.52\% CER from a G2P-only dataset;
Urdu G2P 14.77\% CER (first learned baseline, 4.1$\times$ over
rule-based) and diacritization 3.74\% CER from weak supervision. We
distill cross-cutting lessons for phonological NLP: LLMs are unreliable
phonological labelers; evaluation domains and input formats silently
dominate comparisons; ensembling and curricula fail at data-limited
optima; and byte-level models are the right deployment target for
multilingual TS runtimes.
\end{abstract}

\section{Why phonology for transliteration?}
Transliteration maps (ALA-LC, SBL, RTGS, \dots) assume \emph{vocalized}
input. Undiacritized Arabic or Hebrew, or unsegmented Thai, cannot be
transliterated unambiguously. The models described here supply the
vocalization that makes interscript.org's maps deterministic.

\section{System and results}
Common framework: ByT5 seq2seq (byte-level; tokenizer = UTF-8 bytes, so
TS inference needs no vocabulary assets) except Arabic, where a compact
30M char-level encoder won on cost.

\begin{table}[h]
\centering
\small
\begin{tabular}{llll}
\toprule
Language & Task & Result & Published reference \\
\midrule
Arabic & diacritization & \textbf{2.81/1.69\% DER} & Sadeed 7.29/5.26\%, Gemini 3.19/2.38\% (same benchmark+evaluator) \\
Hebrew & diacritization & \textbf{17.46\% DER} (Biblical) & DictaBERT 35.6\% (same test) \\
Thai & G2P & \textbf{2.32\% PER} & baseline 6.37\% (same test) \\
Persian & G2P/HA & \textbf{89.5\%} / 77.3\% (SB) & Homo-GE2PE 76.9\% \\
Persian & diacritization & \textbf{0.52\% CER} & (none existed) \\
Urdu & G2P & \textbf{14.77\% CER} & epitran 60.0\% \\
Urdu & diacritization & \textbf{3.74\% CER} & (none existed) \\
\bottomrule
\end{tabular}
\end{table}

\section{The bottleneck taxonomy}
\begin{itemize}
  \item \textbf{Arabic --- data quality}: cleaning + 28$\times$ scaling
        halved DER; a 30M encoder matches a 1.5B LM.
  \item \textbf{Hebrew --- data domain}: the SOTA model collapses
        cross-domain (9$\times$); mixed-domain training wins; the teamim
        input-format effect makes cantillation copy-through.
  \item \textbf{Thai --- labeled coverage}: 10K dictionary $\to$ 60K
        via deterministic phonemizer labels; all architectural fixes
        failed, data fixed it.
  \item \textbf{Persian --- representation}: one dataset serves two
        tasks (G2P + diacritization); metric representation choices
        (ezafe) shift results by 18 points.
  \item \textbf{Urdu --- supervision existence}: no gold diacritization
        exists; weak IPA-derived labels at 597K scale suffice.
\end{itemize}

\section{Cross-cutting lessons}
\begin{enumerate}
  \item \textbf{LLMs cannot label phonology.} Tone and haraqat
        hallucination mirrors the known diacritization failure; epitran
        labels are free of it and transfer through fine-tuning.
  \item \textbf{Input formats and domains are hidden variables.}
        Teamim-preserving inputs and Biblical-vs-modern test domain each
        move DER by more than most modeling choices.
  \item \textbf{Beam search is a first-class factor} (12 DER points on
        Hebrew); ensembles must be evaluated under beam search to be
        meaningful.
  \item \textbf{Output-vote ensembling hurts seq2seq restoration}:
        per-character voting lowers error counts yet raises
        edit-distance metrics (incoherent splices).
  \item \textbf{Curricula fail at data-limited optima}; their reported
        gains live in pretraining regimes.
  \item \textbf{Byte-level models are deployment-friendly}: ByT5's
        tokenizer is \texttt{TextEncoder} in TypeScript---no vocab
        files, no WASM sentencepiece---ideal for interscript.org's
        browser/Node targets.
\end{enumerate}

\section{Deployment}
Best checkpoints export to ONNX; distribution follows the Interscript
plan (GitHub Releases primary, HuggingFace canonical, jsDelivr edge).
The TS package wraps onnxruntime-web/node with byte-level tokenization.

\section{Reproducibility}
Per-language repositories (\texttt{rababa}, \texttt{rababa-farsi},
\texttt{rababa-urdu}, \texttt{secryst}) contain training pipelines,
evaluation harnesses, \texttt{docs/RESULTS.md} ground truth, and the
five companion papers in \texttt{docs/paper-*/}.

\end{document}
